\documentclass[11pt]{article}
\usepackage[margin=1in]{geometry}
\usepackage[T1]{fontenc}
\usepackage[utf8]{inputenc}
\usepackage{amsmath,amssymb,amsthm}
\usepackage{enumitem}

\title{ICPC World Finals 2023\\F. Tilting Tiles}
\author{}
\date{}

\begin{document}
\maketitle

\section*{Problem Summary}

We are given a starting board and a target board. A tilt moves every tile as far as possible in one of the
four cardinal directions. Tiles with the same letter are indistinguishable. We must decide whether some
sequence of tilts transforms the start into the target.

\section*{Key Observations}

\begin{itemize}[leftmargin=*]
    \item After we have made at least one horizontal and at least one vertical tilt, every tile is flush in a
    corner. From then on, the occupied shape never changes; only the corner anchor and the order of the
    tiles inside that shape can change.
    \item Therefore every reachable configuration is of one of two types:
    \begin{itemize}[leftmargin=*]
        \item the start itself or a state reachable in one move;
        \item a corner-anchored state obtained after one horizontal and one vertical move, followed by
        repeated cycling through the four corners.
    \end{itemize}
    \item For a fixed corner-anchored base state, one full 4-move cycle induces a permutation $p$ on the
    tile positions. Reachability becomes: does some power $p^n$ produce the target labeling?
    \item On each cycle of $p$, this is a string-rotation problem. Either there is no valid rotation, or the
    valid exponents satisfy one congruence
    \[
    n \equiv a \pmod m.
    \]
    \item The whole board is reachable iff the congruences from all permutation cycles are mutually
    consistent.
\end{itemize}

\section*{Algorithm}

\begin{enumerate}[leftmargin=*]
    \item Compare the multisets of letters in the two boards. If they differ, output \texttt{no}.
    \item Check the trivial cases: the start itself and the four boards reachable in one tilt.
    \item If the board has height $1$ or width $1$, no longer sequence can change anything beyond those
    cases, so we stop.
    \item Otherwise try the eight two-move seeds consisting of one horizontal and one vertical move. Each
    seed produces a corner-anchored base state.
    \item For that base state:
    \begin{itemize}[leftmargin=*]
        \item label the tiles distinctly and simulate one full 4-move corner cycle to obtain the permutation
        $p$ of positions;
        \item for each of the four corner anchors along this cycle, check whether the target has the same
        occupied shape;
        \item if so, read the target letters in the order of the current labels and test whether some power of
        $p$ maps the base letter sequence to that target sequence.
    \end{itemize}
    \item To test a power of $p$, decompose $p$ into cycles. For each cycle:
    \begin{itemize}[leftmargin=*]
        \item extract the source and target strings on that cycle;
        \item find a matching rotation by string matching on the doubled source string;
        \item reduce all valid shifts modulo the minimal period of that cycle string.
    \end{itemize}
    \item If the resulting congruences are pairwise consistent, output \texttt{yes}. If all seeds fail, output
    \texttt{no}.
\end{enumerate}

\section*{Correctness Proof}

We prove that the algorithm returns the correct answer.

\paragraph{Lemma 1.}
After at least one horizontal and at least one vertical tilt, the occupied shape is fixed up to which corner
it is anchored in.

\paragraph{Proof.}
A horizontal tilt compresses each row against one side, and a vertical tilt compresses each column against
one side. After both directions have been used, every row and every column is already maximally
compressed, so later tilts can only move the compressed shape between corners. The set of occupied cells
inside that shape cannot change anymore. \qed

\paragraph{Lemma 2.}
For a fixed corner-anchored base state, the reachable labelings are exactly the boards obtained by
applying some power of the induced permutation $p$ and then possibly stopping at one of the four corner
anchors on that cycle.

\paragraph{Proof.}
By Lemma 1, once the shape is corner-anchored, every later move only changes the corner and the order
of the labeled tiles within the fixed shape. One full round through the four corners applies a fixed
permutation $p$ to the labels. Therefore every later state is described by a power of $p$ together with one
of the four anchors where we stop, and every such state is realized by following that many full rounds and
then the corresponding prefix of the 4-cycle. \qed

\paragraph{Lemma 3.}
For one cycle of the permutation $p$, the set of valid exponents is either empty or one congruence class
$n \equiv a \pmod m$, where $m$ is the minimal period of the source string on that cycle.

\paragraph{Proof.}
Along one cycle of length $L$, applying $p$ repeatedly rotates the letters cyclically. A target labeling is
reachable on that cycle exactly when it is some rotation of the source string. String matching on the
doubled source string finds one such rotation, if it exists.

If the source string has minimal period $m$, then two rotations produce the same string exactly when
their offsets are congruent modulo $m$. Hence all valid exponents form one residue class modulo $m$.
\qed

\paragraph{Theorem.}
The algorithm outputs \texttt{yes} exactly when the target board is reachable.

\paragraph{Proof.}
By Lemma 1, every nontrivial reachable board is covered by one of the eight two-move seeds and one of
the four anchors examined by the algorithm. By Lemma 2, after fixing such a seed and anchor, reachability
reduces exactly to whether some power of the associated permutation $p$ produces the target labeling.

By Lemma 3, each cycle of $p$ contributes either impossibility or one congruence on the exponent. A
single exponent exists exactly when all these congruences are mutually consistent. Therefore the algorithm
accepts exactly the reachable targets. \qed

\section*{Complexity Analysis}

Let $N=hw$ be the number of cells and let $T$ be the number of tiles. Each tilt simulation is $O(N)$. We
try only a constant number of seeds and anchors, so the board-processing cost is $O(N)$. For each
candidate we also process the permutation cycles and run linear-time string matching on the cycle
strings, which is $O(T)$. Thus the total running time is $O(N)$ and the memory usage is $O(N)$.

\section*{Implementation Notes}

\begin{itemize}[leftmargin=*]
    \item Equal-colored tiles are handled by comparing label sequences only after the occupied shape is
    matched.
    \item Congruences are checked by grouping equal moduli and then applying the gcd consistency test.
\end{itemize}

\end{document}
